\documentclass[12pt,a4paper]{article}
\usepackage[a4paper,margin=3cm]{geometry}
\usepackage{xepersian}
\settextfont{Vazirmatn}
\setlatintextfont{DejaVu Serif}
\pagestyle{empty}
\begin{document}
\vspace*{4cm}
\begin{center}
{\LARGE\bfseries وارونی زمان و معمای هستی‌شناختیِ معکوس‌پذیری}

\vspace{0.4cm}
{\large آیا می‌توان «شدن» را معکوس کرد؟}

\vspace{2cm}
{\large ناصر آهنی}

\vspace{0.3cm}
دانشکده فیزیک، دانشگاه صنعتی امیرکبیر، تهران، ایران

\vspace{0.2cm}
{\small \lr{Department of Physics, Amirkabir University of Technology, Tehran, Iran}}

\vspace{0.2cm}
{\small \lr{\texttt{naser.ahani@aut.ac.ir}}}

\vspace{2.5cm}
\rule{6cm}{0.4pt}

\vspace{0.6cm}
\begin{minipage}{0.8\textwidth}
\centering
مقاله‌ی مرجع: «وارونی زمان» (Time Reversal)،\\ نوشته‌ی برایان دبلیو.\ رابرتز

\vspace{0.8cm}
{\small با سپاس از سمینارهای فلسفه فیزیک دانشگاه صنعتی شریف}
\end{minipage}
\end{center}
\end{document}
